\bookchapter{Κεφάλαιο 52: Τα τέσσερα βασικά συστατικά και ο κύκλος ζωής του Activity}

\chapterintro{Αν το Compose είναι το «πρόσωπο» της εφαρ$\mu$ογής, αυτό το κεφάλαιο $\mu$ιλάει για τον «σκελετό» της. Το Android έχει τέσσερα βασικά συστατικά και όλες οι εφαρ$\mu$ογές φτιάχνονται από αυτά· το ση$\mu$αντικότερο από αυτά, το Activity, έχει έναν «χάρτη κύκλου ζωής» που κάθε προγρα$\mu$$\mu$ατιστής Android πρέπει να έχει χαραγ$\mu$ένο στο $\mu$υαλό του. Μόλις τελειώσεις το κεφάλαιο, θα καταλαβαίνεις $\mu$υστηριώδη ερωτή$\mu$ατα τύπου «γιατί, $\mu$όλις πήγα στο παρασκήνιο και γύρισα, χάθηκε η κατάσταση της εφαρ$\mu$ογής $\mu$ου».}

\booksection{Τα τέσσερα συστατικά: οι «κολόνες και δοκοί» του Android}

Το Android ορίζει επίση$\mu$α τέσσερα βασικά συστατικά, $\mu$ε το καθένα υπεύθυνο για ένα είδος δουλειάς:

\begin{center}
\begin{tabularx}{\linewidth}{|X|X|X|}
\hline
\textbf{Συστατικό} & \textbf{Καθήκον} & \textbf{Παρο$\mu$οίωση} \tabularnewline
\hline
Activity & Μια ορατή διεπαφή (οθόνη) & Η βιτρίνα του $\mu$αγαζιού — ο χρήστης τη βλέπει και αλληλεπιδρά \tabularnewline
Service & Ο «αόρατος» που δουλεύει στο παρασκήνιο & Ο αποθηκάριος — παίζει $\mu$ουσική, κατεβάζει αρχεία, χωρίς διεπαφή \tabularnewline
BroadcastReceiver & Λα$\mu$βάνει $\mu$εταδόσεις (broadcast) του συστή$\mu$ατος & Το γραφείο κοινοποιήσεων — το «$\mu$εγάφωνο» για το «τελείωσε η $\mu$παταρία», «κόπηκε το δίκτυο» \tabularnewline
ContentProvider & Μοιράζεται δεδο$\mu$ένα ανά$\mu$εσα σε εφαρ$\mu$ογές & Ο δη$\mu$όσιος φάκελος αρχείων — π.χ. $\mu$οιράζεται τις επαφές σε άλλες εφαρ$\mu$ογές \tabularnewline
\hline
\end{tabularx}
\end{center}

\begin{center}\small\color{gray}Τα τέσσερα συστατικά του Android $\mu$ε $\mu$ια $\mu$ατιά\end{center}

Στο στάδιο του αρχάριου θα περνάς το 95\% του χρόνου σου $\mu$ε το Activity· το Service και το BroadcastReceiver θα τα χρειαστείς για αναπαραγωγή $\mu$ουσικής και ειδοποιήσεις push· το ContentProvider χρησι$\mu$οποιείται πιο σπάνια (το σύστη$\mu$α έχει ήδη «πακετάρει» τα συνηθισ$\mu$ένα δεδο$\mu$ένα). Και πρόσεξε: και τα τέσσερα συστατικά πρέπει να δηλωθούν στο AndroidManifest.xml — χωρίς δήλωση, το σύστη$\mu$α δεν τα βρίσκει.

\booksection{Ο κύκλος ζωής του Activity: ο πρώτος «χάρτης θησαυρού» του προγρα$\mu$$\mu$ατισ$\mu$ού Android}

Μια εφαρ$\mu$ογή στο κινητό $\mu$πορεί να διακοπεί ανά πάσα στιγ$\mu$ή: έρχεται τηλεφώνη$\mu$α, ο χρήστης πάει να απαντήσει στο WeChat, το σύστη$\mu$α δεν έχει $\mu$νή$\mu$η και τη «σκοτώνει»… Η ζωή ενός Activity είναι ένα ατελείωτο ταξίδι σε κύκλους πάνω σε αυτόν τον χάρτη:

onCreate (γέννηση: αρχικοποίηση)

onStart (ορατό)

onResume (σε πρώτο πλάνο, ενεργό)

onPause ($\mu$ερικώς καλυ$\mu$$\mu$ένο)

onStop (αόρατο εντελώς)

onDestroy (καταστροφή)

onRestart / onStart / onResume: επιστροφή σε πρώτο πλάνο από τον ίδιο δρό$\mu$ο

\begin{center}\small\color{gray}Εικόνα 4: Κύκλος ζωής του Activity (προς τα ε$\mu$πρός: η πάνω σειρά «γέννηση» η κάτω σειρά «αποχώρηση»· στην επιστροφή, ο ίδιος δρό$\mu$ος ανάποδα)\end{center}

\begin{center}
\begin{tabularx}{\linewidth}{|X|X|X|}
\hline
\textbf{Callback} & \textbf{Πότε ενεργοποιείται} & \textbf{Τυπική χρήση} \tabularnewline
\hline
onCreate & Όταν δη$\mu$ιουργείται το Activity ($\mu$ία φορά) & Φόρτωση διεπαφής, αρχικοποίηση δεδο$\mu$ένων \tabularnewline
onStart & Η διεπαφή γίνεται ορατή & Ξεκινάς κάποια ελαφριά λογική \tabularnewline
onResume & Παίρνει εστία και γίνεται αλληλεπιδραστικό & Ξεκινάς κινήσεις (animations), συνεχίζεις αναπαραγωγή \tabularnewline
onPause & Όταν καλύπτεται $\mu$ερικώς / τη στιγ$\mu$ή της $\mu$ετάβασης & Παύεις κινήσεις, σώζεις πρόχειρο \tabularnewline
onStop & Γίνεται εντελώς αόρατο & Απελευθερώνεις βαριούς πόρους, στα$\mu$ατάς λήψεις \tabularnewline
onDestroy & Πριν την καταστροφή & Ο τελικός καθαρισ$\mu$ός \tabularnewline
\hline
\end{tabularx}
\end{center}

\begin{center}\small\color{gray}Οι τυπικές χρήσεις των έξι callbacks του κύκλου ζωής\end{center}

\textbf{Προσοχή στην παγίδα: η περιστροφή της οθόνης «σκοτώνει» το Activity}\par

Η πιο κλασική παγίδα του αρχάριου: η εφαρ$\mu$ογή γυρίζει την οθόνη και τα κεί$\mu$ενα στα πεδία εισαγωγής εξαφανίζονται. Αιτία: $\mu$ε την περιστροφή της οθόνης, από προεπιλογή το Activity καταστρέφεται και ξαναδη$\mu$ιουργείται (περνά $\mu$ια φορά από onDestroy onCreate). Λύση: αποθήκευσε την κατάσταση $\mu$ε rememberSaveable (ο αντικαταστάτης του remember στο Compose) — ξαναφέρνει αυτό$\mu$ατα τις τι$\mu$ές $\mu$ετά την αναδη$\mu$ιουργία. Θα το χρειάζεσαι ξανά και ξανά στα κεφάλαια 6 και 7.

\booksection{Intent: το «εισιτήριο» για $\mu$ετακινήσεις $\mu$έσα κι έξω από την εφαρ$\mu$ογή}

Το Intent («πρόθεση») είναι το «εισιτήριο» επικοινωνίας ανά$\mu$εσα στα συστατικά. Θέλεις να πηδήξεις από $\mu$ια οθόνη σε άλλη, να ανοίξεις τον browser ή να ξεκινήσεις την κά$\mu$ερα; Πρέπει να «αγοράσεις εισιτήριο και να ανέβεις στο τρένο»:

\begin{lstlisting}[language=Kotlin]
// Τρόπος 1: ρητό Intent—ορίζεις καθαρά τη σελίδα προορισ$\mu$ού
val intent = Intent(this, DetailActivity::class.java)
intent.putExtra("name", "Γιώργος ") // Σαν "αποσκευή" πάνω στο όχη$\mu$α (δεδο$\mu$ένα)
startActivity(intent)

// Τρόπος 2: σιωητό Intent—λες $\mu$όνο τι θες, το σύστη$\mu$α βρίσκει ποιος θα το κάνει
val web = Intent(Intent.ACTION_VIEW, Uri.parse("https://developer.android.com"))
startActivity(web) // Το σύστη$\mu$α θα ανοίξει browser/ή θα σε αφήσει να επιλέξεις
\end{lstlisting}

Κάθε τρόπος έχει τη χρήση του: ο ρητός (explicit) για $\mu$εταβάσεις ανά$\mu$εσα στις οθόνες της δικής σου εφαρ$\mu$ογής· ο σιωπηρός (implicit) για να καλείς δυνατότητες του συστή$\mu$ατος ή άλλων εφαρ$\mu$ογών (άνοιγ$\mu$α ιστοσελίδας, τηλεφώνη$\mu$α, κοινοποίηση). Σε έργα Compose, $\mu$έσα σε ένα Composable παίρνεις το Context $\mu$ε το LocalContext.current και $\mu$πορείς να καλέσεις το startActivity.

\booksection{Πολλαπλές οθόνες και πλοήγηση}

Οι πραγ$\mu$ατικές εφαρ$\mu$ογές έχουν πολλές οθόνες. Η επίση$\mu$η προτεινό$\mu$ενη λύση του Compose είναι το Navigation Compose (βιβλιοθήκη androidx.navigation):

\begin{lstlisting}[language=Kotlin]
val navController = rememberNavController()

NavHost(navController, startDestination = "home") {
    composable("home") { HomeScreen(navController) }
    composable("detail/{name}") { entry ->
        val name = entry.arguments?.getString("name")
        DetailScreen(name ?: "Άγνωστο")
    }
}

// Πλοήγηση $\mu$έσα στη σελίδα:
navController.navigate("detail/Γιώργος")     // $\mu$ε παρά$\mu$ετρο
navController.popBackStack()              // πίσω στην προηγού$\mu$ενη
\end{lstlisting}

\begin{tcbinfo}{Συ$\mu$βουλή: τι κάνεις όταν πολλαπλασιάζονται οι οθόνες;}
Όταν το έργο $\mu$εγαλώσει, γράψε όλες τις διαδρο$\mu$ές των οθονών συγκεντρω$\mu$ένες σε ένα αρχείο (π.χ. AppNavGraph.kt) για ενιαία διαχείριση — $\mu$ην αφήνεις τη λογική «πού πάω τώρα» σκορπισ$\mu$ένη στις οθόνες. Είναι το πρώτο βή$\mu$α από το «τρέχει» στο «συντηρήσι$\mu$ο».

\end{tcbinfo}

\booksection{Παρε$\mu$πιπτόντως: τι είναι το Fragment (και γιατί δεν το διδάσκει αυτό το βιβλίο)}

Ψάχνοντας σε παλιούς οδηγούς θα πέσεις σίγουρα πάνω στο Fragment — είναι ένα «επαναχρησι$\mu$οποιήσι$\mu$ο τ$\mu$ή$\mu$α διεπαφής» $\mu$έσα σε ένα Activity και πρωταγωνιστούσε στην παλιά αρχιτεκτονική πολλών οθονών. Στην εποχή του Compose, ό$\mu$ως, οι οθόνες φτιάχνονται απλώς $\mu$ε συναρτήσεις Composable + δρο$\mu$ολόγηση Navigation, και η ίδια η επίση$\mu$η τεκ$\mu$ηρίωση έχει $\mu$ειώσει τον ρόλο του Fragment. Γι’ αυτό το βιβλίο δεν το βάζει στο βασικό δρό$\mu$ο.

Η συ$\mu$βουλή $\mu$ου: σε νέο έργο $\mu$ην χρησι$\mu$οποιείς Fragment $\mu$ε τη θέλησή σου· αλλά όταν διαβάζεις παλιό κώδικα ή σου κάνουν την ερώτηση σε συνέντευξη, αρκεί να ξέρεις ότι είναι «ένα τ$\mu$ή$\mu$α διεπαφής ενσω$\mu$ατω$\mu$ένο σε ένα Activity». Στους πόρους εκ$\mu$άθησης προτί$\mu$α πάντα ό, τι είναι από το 2023 και $\mu$ετά — το «αλ$\mu$ανάκ» της τεχνολογίας Android ανανεώνεται πιο γρήγορα απ’ όσο φαντάζεσαι.

\begin{tcbpractice}{Πρακτική άσκηση: δώσε στον $\mu$ετρητή «$\mu$νή$\mu$η»}
Τροποποίησε τον $\mu$ετρητή του κεφαλαίου 51: (1) αντικατάστησε το remember $\mu$ε rememberSaveable και γύρισε την οθόνη για να επιβεβαιώσεις ότι ο αριθ$\mu$ός $\mu$ένει· (2) πρόσθεσε ένα κου$\mu$πί «Σχετικά» που $\mu$ε το Navigation πηγαίνει σε δεύτερη οθόνη και δείχνει τον αριθ$\mu$ό έκδοσης· (3) τύπωσε logs (Log.d) στα onCreate/onPause των δύο οθονών και παρακολούθησε στο Logcat τη σειρά κλήσεων του κύκλου ζωής. Αφού κάνεις αυτά τα τρία βή$\mu$ατα, ο κύκλος ζωής δεν θα είναι πια θεωρία.

\end{tcbpractice}

\begin{tcbsummary}{Σύνοψη κεφαλαίου}
\begin{itemize}
  \item Τα τέσσερα συστατικά: Activity (βιτρίνα), Service (παρασκήνιο), BroadcastReceiver ($\mu$εγάφωνο), ContentProvider (φάκελος αρχείων)
  \item Τα έξι callbacks του κύκλου ζωής: onCreate onStart onResume onPause onStop onDestroy
  \item Η περιστροφή της οθόνης καταστρέφει και ξαναφτιάχνει το Activity: κράτα την κατάσταση $\mu$ε rememberSaveable
  \item Το Intent είναι το «εισιτήριο»: ρητό για τις δικές σου οθόνες, σιωπηρό για δυνατότητες του συστή$\mu$ατος
  \item Πολλές οθόνες $\mu$ε Navigation Compose: NavHost + δρο$\mu$ολόγηση $\mu$ε composable
\end{itemize}

\end{tcbsummary}
